% ============================== CI-1 ==============================
\reading{CI-1}{Advances in Artificial Intelligence: Implications for Capital Markets Activities}{STRIKE FIRST}{6}

\begin{cikeybox}[title={The four objectives, in the spine's own words}]
\begin{enumerate}[label=\textbf{\alph*.}]
\item Describe current uses of artificial intelligence (AI) and machine learning (ML) in
capital markets, and potential future uses of sophisticated AI models including GenAI.
\item Explain the implications of further adoption of AI on market dynamics.
\item Explain the implications of further adoption of AI on financial stability.
\item Describe best practices for regulators pertaining to AI, including the use of AI by
supervisors and recommendations for policy on the use of AI by regulated entities.
\end{enumerate}
\greyline{Source: ``Advances in Artificial Intelligence: Implications for Capital Markets
Activities'', IMF (October 2024). The class notes label this Reading 100; it is CI-1 on the
2026 spine.}
\end{cikeybox}

% ------------------------------------------------------------------
\lo{CI-1 a}{Describe current uses of artificial intelligence (AI) and machine learning (ML) in capital markets, and potential future uses of sophisticated AI models including GenAI.}

\subsection*{Overview of AI and ML in capital markets}

This chapter is based on an IMF survey of 27 firms across buy-side and sell-side
participants.

AI models are classified into two groups:

\begin{enumerate}
\item \term{Traditional AI/ML}: NLP, decision trees, clustering.
\item \term{Advanced AI}: LLMs, deep learning, reinforcement learning.
\end{enumerate}

\begin{cifigblock}
{\centering
\renewcommand{\arraystretch}{1.25}
\begin{tabularx}{\textwidth}{@{}>{\raggedright\arraybackslash}X >{\raggedright\arraybackslash}X@{}}
\toprule
\rowcolor{PaleNavy}
{\sffamily\bfseries\color{FRMNavy}Key benefits} &
{\sffamily\bfseries\color{FRMRed}Key risks}\\
\midrule
Higher efficiency, productivity, and forecasting accuracy &
Operational risks (model errors, system failures)\\
Improved investment frameworks and crash risk measurement &
Increased volatility from rapid trading under stress\\
Better use of unstructured and alternative data &
Cybersecurity and market manipulation threats\\
 & Lack of transparency and model opacity\\
\bottomrule
\end{tabularx}
\par}
\figcap{Table CI-1.1. The benefit column and the risk column are not independent lists.
Each risk is the shadow of the benefit beside it: speed buys efficiency and costs stability,
and the use of unstructured data buys insight and costs transparency.}
\end{cifigblock}

\subsection*{Key applications in capital markets}

\begin{cifigblock}
{\centering
\begin{tikzpicture}[
  font=\sffamily\scriptsize,
  proc/.style={draw=FRMRule, fill=white, text width=2.05cm, align=center,
               minimum height=0.95cm, inner sep=2pt},
  sub/.style={draw=FRMRule, fill=PaleNavy, text width=2.05cm, align=center,
              minimum height=1.05cm, inner sep=2pt},
  hdr/.style={draw=FRMRule, fill=FRMSoft, text width=1.5cm, align=center,
              minimum height=2.1cm, inner sep=2pt}
]
\node[hdr] (key) at (0,0) {\bfseries\color{FRMNavy}Key\\processes};

\node[proc] (p1) at (2.05,0.55) {Client / institution profiling};
\node[proc, text width=4.35cm] (p2) at (5.25,0.55) {Asset allocation};
\node[proc] (p3) at (8.45,0.55) {Trading};
\node[proc] (p4) at (10.7,0.55) {Risk management};

\node[sub] (s1) at (2.05,-0.55) {Identification of needs and constraints};
\node[sub, text width=1.35cm] (s2) at (3.65,-0.55) {Asset class allocation};
\node[sub, text width=1.35cm] (s3) at (5.25,-0.55) {Sectoral allocation};
\node[sub, text width=1.35cm] (s4) at (6.85,-0.55) {Security selection};
\node[sub] (s5) at (8.45,-0.55) {Orders placement and execution};
\node[sub, text width=1.35cm] (s6) at (10.15,-0.55) {Risk monitoring};
\node[sub, text width=1.35cm] (s7) at (11.5,-0.55) {Reporting};
\end{tikzpicture}
\par}
\figcap{Figure CI-1.1. Where AI actually sits in the capital markets workflow. Read the top
row as the four processes and the bottom row as the steps inside them: asset allocation is
the widest because it decomposes into three separate decisions, which is why it absorbs the
most AI effort. Sources given in the notes: academic studies; IMF outreach discussions
(see Box 3.1); prospectuses from third-party services; IMF staff compilations.}
\end{cifigblock}

\begin{cidefbox}[title={The point to hold on to}]
AI in finance is not new. Tools like \term{robo-advisors} and \term{AI-driven ETFs} already
exist. While robo-advisory AUM has grown rapidly, AI-powered ETFs still make up a small
portion of the market.
\end{cidefbox}

\subsection*{Role of generative AI}

\begin{itemize}
\item Generative AI (GenAI) is building on existing models by improving investment idea
generation, coding efficiency, and client-facing processes. LLMs are increasingly used for
forecasting and analytics.
\item While AI can potentially lower entry barriers and speed up market reactions, a
\term{human-in-the-loop approach remains critical} to manage regulatory and ethical risks.
\end{itemize}

\subsection*{Emerging trends and future applications}

AI adoption is strongest in liquid markets such as equities, FX, and bonds, where systems
can learn quickly from dynamic data. At the same time, high infrastructure and talent costs
may lead to market concentration, favouring large data-rich firms.

AI applications span a broad range of use cases, including the development of
forward-looking indicators, integration of alternative datasets, extraction of trading
signals, portfolio allocation and optimisation, pricing and forecasting, risk assessment,
and trade automation.

\begin{cifigblock}
{\centering
\renewcommand{\arraystretch}{1.2}
\begin{tabularx}{\textwidth}{@{}>{\raggedright\arraybackslash}X
                              >{\raggedright\arraybackslash}X
                              >{\raggedright\arraybackslash}X@{}}
\toprule
\rowcolor{PaleGreen}
{\sffamily\bfseries\color{FRMGreen}Benefits of AI} &
{\sffamily\bfseries\color{FRMNavy}AI in asset management} &
{\sffamily\bfseries\color{FRMGold}Regulation benefits of AI}\\
\midrule
Efficiency goes up &
AI supports decisions, it does not replace them &
SupTech: used by regulators (monitoring, anomaly detection)\\
Better data processing, especially unstructured data &
Humans still take final calls &
RegTech: used by firms (compliance, reporting)\\
Lower trading costs &
Fully autonomous trading still limited &
Helps improve surveillance and risk oversight\\
Narrower bid--ask spreads &
AI mainly feeds into existing models & \\
Improved forecasting accuracy & & \\
\bottomrule
\end{tabularx}
\par}
\figcap{Table CI-1.2. The middle column is the one that decides questions. Three of its four
lines say the same thing from different angles: the human stays in the decision seat.}
\end{cifigblock}

AI will increasingly drive \term{investment and trading decisions}, with higher autonomy,
especially in liquid markets. However, a \term{human-in-the-loop approach remains critical}
for oversight, control, and regulatory compliance.

The growing number of AI-related patents highlights rapid innovation, particularly in
high-frequency trading and asset management functions like rebalancing, forecasting, and
valuation, along with expanding use of alternative data and new asset classes. Overall, AI
adoption will deepen, but decisions will still be guided by humans.

\subsection*{Broader implications and opportunities}

\begin{itemize}
\item Demand for AI talent is rising sharply in finance, especially in roles linked to
trading and investments. AI can support emerging markets by improving liquidity, capital
access, and financial services, with GenAI structuring unorganised data and synthetic data
enabling model training.
\item However, risks such as bias and misrepresentation remain. Despite some job
automation, AI is expected to enhance financial inclusion and reduce barriers to entry
across markets.
\end{itemize}

% ------------------------------------------------------------------
\lo{CI-1 b}{Explain the implications of further adoption of AI on market dynamics.}

\subsection*{AI adoption and shifts in market structure}

AI is reshaping capital markets by favouring faster, more flexible players. \term{NBFIs},
which hold over half of global financial assets, are better positioned to benefit from AI
compared to traditional banks, which face legacy and regulatory constraints.

Algorithmic trading is a key driver, accounting for about 70\% of equities, over 50\% of
futures, and more than 40\% of options trading in the U.S., but less than 10\% in fixed
income, showing uneven adoption across markets.

\begin{cifigblock}
{\centering
\begin{tikzpicture}
\begin{axis}[
  width=11.2cm, height=4.6cm,
  xbar, bar width=12pt,
  xmin=0, xmax=108,
  xlabel={\sffamily\scriptsize Share of U.S. trading activity that is algorithmic (\%)},
  symbolic y coords={Fixed income, Options, Futures, Equities},
  ytick=data,
  tick label style={font=\sffamily\scriptsize},
  label style={font=\sffamily\scriptsize},
  nodes near coords,
  nodes near coords style={font=\sffamily\scriptsize\color{FRMNavy},
     fill=white, inner sep=1.5pt, anchor=west, xshift=3pt},
  point meta=explicit symbolic,
  xmajorgrids, grid style={FRMRule, dashed},
  axis line style={FRMRule},
  enlarge y limits=0.22
]
\addplot[fill=FRMNavy!75, draw=FRMNavy] coordinates {
  (10,{Fixed income}) [under 10\%]
  (40,{Options})      [over 40\%]
  (50,{Futures})      [over 50\%]
  (70,{Equities})     [about 70\%]
};
\end{axis}
\end{tikzpicture}
\par}
\figcap{Figure CI-1.2. Adoption is not uniform, and the ordering is the examinable part:
equities, then futures, then options, then a long drop to fixed income. The fixed-income bar
is drawn at its ceiling, since the notes give it only as ``less than 10\%''. The gap between
the third and fourth bars is what ``uneven adoption across markets'' means.}
\end{cifigblock}

Overall, high AI costs are increasing market concentration, giving larger firms a strong
advantage while smaller firms rely on outsourcing, leading to dominance by a few major
players.

\subsection*{Algorithmic trading and market implications}

\begin{enumerate}
\item \term{Market efficiency and liquidity.} Algorithmic trading improves efficiency and
liquidity in normal conditions by processing large volumes of data. However, during stress,
liquidity can decline as informed traders and high-frequency traders withdraw, leading to
higher volatility.
\item \term{Price behaviour.} It has reduced unnecessary price fluctuations in stable
markets, but can contribute to sharp price swings in illiquid or stressed environments.
\item \term{Market stability and risk limits.} Risk controls such as automatic shutdowns and
trading limits, while designed to protect markets, can sometimes reduce stability during
stress by triggering feedback loops and draining liquidity. This leads to ``flighty
liquidity under stress'', where high-frequency traders step back, order cancellations
decline, and institutional investors increasingly rely on hidden orders to minimise market
impact.
\item \term{Impact of GenAI.} The integration of GenAI can further transform algorithmic
trading by lowering barriers to entry through improved coding, testing, and automation. It
also enables better processing of complex text-based data and supports more standardised and
scalable trading and risk management across different asset classes.
\end{enumerate}

\begin{citrapbox}[title={Direction matters in all four of the points above}]
Every one of them is the same shape: a benefit in \emph{normal} conditions and the reverse
effect in \emph{stressed} conditions. Efficiency and liquidity improve normally and decline
under stress; price fluctuations fall in stable markets and swing sharply in illiquid ones;
risk controls protect markets normally and can drain liquidity under stress. The notes state
the condition every time, so the condition is the discriminating word.
\end{citrapbox}

\begin{cifigblock}
{\centering
\renewcommand{\arraystretch}{1.2}
\small
\begin{tabularx}{\textwidth}{@{}p{2.3cm}
                              >{\raggedright\arraybackslash}X
                              >{\raggedright\arraybackslash}X@{}}
\toprule
\rowcolor{FRMSoft}
&{\sffamily\bfseries\color{FRMRed}Negative scenario} &
{\sffamily\bfseries\color{FRMGreen}Positive scenario}\\
\midrule
\term{Market liquidity} &
AI magnifies existing risks related to algorithmic trading by facilitating its growth. AI
could ``democratize'' and expand algorithmic trading activity to a broader set of assets and
geographic areas. This could exacerbate risks related to sudden liquidity withdrawal under
stressed conditions. &
AI increases the stability of algorithmic trading under stressed conditions. AI-driven
algorithms could operate in a wider set of market conditions than traditional algorithms,
with lower flash-crash risk, and reduced liquidity withdrawal under stress.\\
\addlinespace
\term{Leverage} &
AI-driven strategies boost short-term leverage. As arbitrage opportunities are exploited
more efficiently by more advanced algorithms, remaining opportunities might require higher
leverage to deliver similar returns. &
AI improves the management of leverage and related risks. AI could facilitate more frequent
and automated management of leveraged positions, based on more inputs, and mitigate
operational lags.\\
\addlinespace
\term{Inter\-connected\-ness} &
\multicolumn{2}{@{}>{\raggedright\arraybackslash}p{\dimexpr\textwidth-2.6cm\relax}@{}}{%
AI increases interconnectedness. AI could proliferate algorithmic trading to other asset
classes, geographic regions, and trading venues, and also operate in between different
market segments; that is, in a multi-asset and multitrading venue approach.}\\
\addlinespace
&
Increased interconnectedness leads to higher correlations between capital market segments,
facilitating spillovers and transmission of stress. &
Market access, efficiency, and liquidity improve for some market segments, including
emerging markets.\\
\bottomrule
\end{tabularx}
\par}
\figcap{Table CI-1.3 (the notes' Table 3.1). Note the structural oddity in the third row,
which is in the source and is kept: interconnectedness \emph{increases} in both scenarios.
Only the consequence differs, so no answer that turns on whether interconnectedness rises
can be resolved from the scenario alone. Source: IMF staff assessment.}
\end{cifigblock}

\subsection*{Emerging market dynamics and financial stability risks}

\begin{enumerate}[label=\textbf{\alph*)}]
\item \term{Faster market reactions to news.} AI enables quicker and more accurate
interpretation of information (e.g., FOMC releases), making initial market responses closer
to final outcomes.
\item \term{Higher and procyclical trading volumes.} AI-driven strategies adjust portfolios
rapidly, leading to higher turnover. This supports price discovery in stable markets but can
amplify volatility during stress.
\item \term{Potential for AI-driven manipulation or collusion.} Interactions between
algorithms may create mispricing or strategic advantages, increasing the risk of
inefficiencies and unintended collusive behaviour.
\end{enumerate}

% ------------------------------------------------------------------
\lo{CI-1 c}{Explain the implications of further adoption of AI on financial stability.}

\subsection*{Impact of AI on financial stability}

Several areas of concern regarding the potential impact of AI on financial stability were
highlighted by participants in the IMF outreach survey. These include the following:

\begin{enumerate}[label=\textbf{\alph*)}]
\item \term{Concentration risk (herding).} A key concern arises when trading strategies rely
heavily on similar open-source AI models and limited data vendors. This creates vendor
concentration, where dependence on a few dominant providers can pose systemic risk if they
fail.
\item \term{Market fragility.} AI-driven markets may face excess volatility, flash crashes,
illiquidity, cyber risks, and data poisoning due to rapid automated decisions and weak
guardrails. There is also a shortage of skilled professionals to manage such systems
effectively.
\item \term{Model hallucinations and lack of explainability.} Unreliable or opaque AI outputs
can reduce trust, making market participants hesitant to rely on automated decisions.
\item \term{Market manipulation.} AI can be misused for misinformation and deepfake content,
threatening market integrity and investor confidence.
\item \term{High fixed costs.} Significant infrastructure and talent costs favour large firms,
reinforcing market concentration and competitive imbalances.
\end{enumerate}

\term{Cyber risks} are particularly significant, including impersonation, data manipulation,
and system disruptions, which can erode confidence and trigger market instability.
Additional concerns include data misuse, fraud, and the transmission of shocks across
economies.

\subsection*{Regulatory responses and systemic implications}

Survey participants expect regulators to provide clearer guidance on AI risk management,
including transparency, disclosures, and extreme-scenario stress testing. There is a strong
emphasis on accountability frameworks and balanced regulation that supports innovation while
ensuring customer protection. Flexible, principles-based guidelines are seen as more
effective than rigid rules.

\subsection*{The four financial stability challenges}

Financial stability challenges arising from the adoption of AI in capital markets include the
following:

\begin{enumerate}
\item \term{Higher volatility and market speed under stress.} While AI improves efficiency,
liquidity, and price discovery, it can amplify volatility during stress due to leverage,
herding behaviour, and similar model outputs. This may trigger rapid selloffs and
flight-to-safety dynamics, with some AI models potentially failing or behaving
inconsistently during extreme events.
\item \term{Monitoring and opacity challenges.} As activities shift from banks to NBFIs,
lower regulatory oversight may reduce transparency. Increased interconnectedness across
markets, asset classes, and regions due to AI-driven strategies can further complicate
monitoring and risk assessment.
\item \term{Cyber risks and market manipulation.} AI increases exposure to cyber threats,
data breaches, and manipulation through misinformation or deepfakes, potentially distorting
asset prices and undermining market trust.
\item \term{Operational risks from third-party dependence.} Heavy reliance on a few AI and IT
service providers creates systemic vulnerability, where disruptions or failures of key
providers could have widespread market impact.
\end{enumerate}

\begin{cifigblock}
{\centering
\begin{tikzpicture}[
  font=\sffamily\scriptsize,
  src/.style={draw=FRMRule, fill=PaleNavy, text width=3.5cm, align=left,
              minimum height=0.82cm, inner sep=4pt, rounded corners=1pt},
  chg/.style={draw=FRMRule, fill=PaleRed, text width=3.9cm, align=left,
              minimum height=0.82cm, inner sep=4pt, rounded corners=1pt},
  arr/.style={-{Stealth[length=4pt]}, FRMGold, line width=0.7pt}
]
\node[font=\sffamily\bfseries\scriptsize\color{FRMNavy}] at (0,1.55)
  {What the survey participants worried about};
\node[font=\sffamily\bfseries\scriptsize\color{FRMRed}] at (7.4,1.55)
  {What it becomes at system level};

\node[src] (a) at (0,1.00)  {a) Concentration / herding};
\node[src] (b) at (0,0.00)  {b) Market fragility};
\node[src] (c) at (0,-1.05) {c) Hallucination, no explainability};
\node[src] (d) at (0,-2.05) {d) Market manipulation};
\node[src] (e) at (0,-3.05) {e) High fixed costs};

\node[chg] (c1) at (7.6,1.00)  {1. Higher volatility and market speed under stress};
\node[chg] (c2) at (7.6,-0.35) {2. Monitoring and opacity challenges};
\node[chg] (c3) at (7.6,-1.70) {3. Cyber risk and market manipulation};
\node[chg] (c4) at (7.6,-3.05) {4. Operational risk from third-party dependence};

\draw[arr] (a.east) -- ([yshift=3pt]c1.west);
\draw[arr] (b.east) -- ([yshift=-3pt]c1.west);
\draw[arr] (c.east) -- ([yshift=3pt]c2.west);
\draw[arr] (e.east) -- ([yshift=-3pt]c2.west);
\draw[arr] (d.east) -- (c3.west);
\draw[arr] (a.east) -- ([yshift=3pt]c4.west);
\draw[arr] (e.east) -- ([yshift=-3pt]c4.west);

\node[fill=white, inner sep=1.5pt, text width=2.6cm, align=center,
      font=\sffamily\scriptsize\itshape\color{FRMGrey}] at (3.8,-4.05)
  {Concentration and cost each feed two channels};
\end{tikzpicture}
\par}
\figcap{Figure CI-1.3. The five concerns on the left and the four stability challenges on the
right are not two lists, they are the same material at firm level and at system level.
Concentration (a) and high fixed costs (e) each feed two channels, which is why they are the
hardest to regulate away: they show up as volatility and as third-party operational risk at
the same time.}
\end{cifigblock}

% ------------------------------------------------------------------
\lo{CI-1 d}{Describe best practices for regulators pertaining to AI, including the use of AI by supervisors and recommendations for policy on the use of AI by regulated entities.}

\subsection*{Regulatory and supervisory scope}

Existing capital market regulations also apply to AI systems, but institutions remain fully
responsible for their use, whether developed internally or sourced from third parties.

\subsection*{Best practices and supervisory use of AI}

\begin{enumerate}[label=\textbf{\alph*)}]
\item \term{Regulators are expected to stay proactive as AI adoption accelerates.}
Supervisors must build capabilities to process large datasets and monitor complex
strategies, while ensuring frameworks remain adaptable. Key focus areas include governance,
transparency and explainability, outsourcing risks, data quality and bias, reporting
standards, and maintaining human oversight.
\item \term{Role of global regulators.} Bodies such as the Financial Stability Board, Basel
Committee on Banking Supervision, Bank for International Settlements, National Institute of
Standards and Technology, and International Organization of Securities Commissions are
strengthening frameworks to enhance investor protection, market integrity, and financial
stability. This includes initiatives like NIST's AI framework and IOSCO's focus on
algorithmic trading risks and market intermediation.
\item \term{Policy direction and recommendations.} Safe AI adoption requires evaluating
current regulations, developing new risk management approaches, and promoting governance
standards through public--private collaboration. Frameworks must address robustness,
explainability, cybersecurity, and data privacy concerns.
\item \term{Supervisory benefits of AI.} AI enables efficiency gains through automated data
validation, integration of diverse datasets, and detection of abnormal trading patterns. It
also supports real-time monitoring and identification of misinformation. Advanced AI systems
further enhance capabilities such as code optimisation, information retrieval, and analytical
support.
\end{enumerate}

\subsection*{Policy recommendations and regulatory priorities}

Given the financial stability risks faced by both the banking and NBFI sectors, regulatory
and supervisory frameworks for AI should be strengthened to balance the realisation of
benefits with the need for enhanced risk management. The following areas have been
identified as \term{key priorities to address}:

\begin{enumerate}
\item \term{Market speed and volatility under stress.} Faster market speeds can increase
volatility, especially during stress. Regulators may need to redesign or enhance tools like
circuit breakers, test trading algorithms in both normal and stressed conditions, and
reassess margin frameworks, including initial margins, margin call volatility, variation
margins, and the responsiveness of central counterparties' models during crises.
\item \term{Operational risk from third-party dependence.} Heavy reliance on a few AI and IT
providers creates systemic vulnerability. Regulators should map interconnections across
providers and strengthen cyber resilience, ensuring systems can prevent, withstand, and
recover from attacks, especially those targeting AI models and data.
\item \term{Opacity and monitoring challenges.} Financial institutions must map
interdependencies between AI models, data sources, and infrastructure. Risks arise when data
is incomplete (e.g., not covering full financial cycles) or concentrated among a few
providers. Regulators should require NBFIs to disclose AI usage, while also strengthening
liquidity buffers and risk management to reduce the risk of liquidity runs and mispricing,
especially given the growing importance of large traders in AI-driven markets.
\item \term{OTC market resilience.} OTC markets need stronger resilience and oversight as AI
adoption expands. Regulators may require more detailed transaction reporting and consider
policies such as improved liquidity frameworks, greater use of central clearing, stronger
market-making support, and margin requirements for non-cleared derivatives. Additionally,
liquidity and funding backstops should be developed to stabilise markets during extreme
stress.
\end{enumerate}

\begin{cikeybox}[title={Two actors, two verbs}]
This objective splits cleanly in half and the notes keep the halves apart. \emph{Supervisors
using AI} is item (d) above: validation, integration, anomaly detection, real-time
monitoring. \emph{Policy for regulated entities using AI} is the numbered priority list:
circuit breakers and margin frameworks, third-party mapping, NBFI disclosure, OTC resilience.
The first is a tool, the second is a rule.
\end{cikeybox}
